% Generated by scripts/make_numbers.py from results/benchmark/*.csv -- do not edit by hand.
\newcommand{\RzeroKL}{99.1}  % R0 membrane+bending, degree 4: non-affine energy captured [%]
\newcommand{\RzeroKLDims}{23}  % R0 membrane+bending degree 4: dimensions
\newcommand{\RzeroGeneric}{96.2}  % R0 generic polynomial, degree 4 [%]
\newcommand{\RzeroGenericDims}{24}  % R0 generic polynomial degree 4: dimensions
\newcommand{\RzeroMembrane}{45.6}  % R0 membrane only, degree 4 [%]
\newcommand{\RzeroKLResidNm}{0.58}  % R0 membrane+bending: median RMS left [nm]
\newcommand{\RzeroGenericResidNm}{1.18}  % R0 generic polynomial: median RMS left [nm]
\newcommand{\NdmaxOurs}{64}  % E1: FEM runs to reach median R2(d_max) >= 0.95, ours
\newcommand{\NdmaxBaseline}{512}  % E1: same, best learned baseline
\newcommand{\EoneOursThirtyTwo}{0.925}  % E1: R2(d_max) at N=32, ours
\newcommand{\EoneBaseThirtyTwo}{0.791}  % E1: R2(d_max) at N=32, best learned baseline
\newcommand{\EoneMembraneDmaxSixtyFour}{0.959}  % E1: R2(d_max) at N=64, membrane-only basis
\newcommand{\EoneOursDmaxSixtyFour}{0.956}  % E1: R2(d_max) at N=64, ours
\newcommand{\EoneNoPriorDmaxSixtyFour}{0.887}  % E1: R2(d_max) at N=64, ours without analytic prior
\newcommand{\EoneOursResidOneTwoEight}{0.838}  % E1: R2(resid_3sigma) at N=128, ours
\newcommand{\EoneMembraneResidOneTwoEight}{0.670}  % E1: R2(resid_3sigma) at N=128, membrane-only
\newcommand{\EoneOursFieldFiveTwelve}{15.5}  % E1: field RMSE at N=512, ours [nm]
\newcommand{\EoneMembraneFieldFiveTwelve}{18.5}  % E1: field RMSE at N=512, membrane-only [nm]
\newcommand{\EtwoOursFour}{2.5}  % E2: 4 marks, 2 nm noise, 256-run prior: median RMSE [nm]
\newcommand{\EtwoOursFourSmallPrior}{4.4}  % E2: same with 64-run prior [nm]
\newcommand{\EtwoAffineFour}{8.1}  % E2: 4 marks, 2 nm noise, affine overlay model [nm]
\newcommand{\EtwoPriorOnly}{8.1}  % E2: 256-run prior, no marks [nm]
\newcommand{\EtwoPriorAffineFour}{2.6}  % E2: 4 marks, 2 nm noise, 256-run prior + affine fit [nm]
\newcommand{\EtwoOursNine}{2.5}  % E2: 9 marks, 2 nm noise, 256-run prior, ours [nm]
\newcommand{\EtwoPriorAffineNine}{2.4}  % E2: 9 marks, 2 nm noise, 256-run prior + affine fit [nm]
\newcommand{\EtwoOursSixteen}{1.4}  % E2: 16 marks, 2 nm noise, 256-run prior, ours [nm]
\newcommand{\EtwoPriorAffineSixteen}{2.2}  % E2: 16 marks, 2 nm noise, 256-run prior + affine fit [nm]
\newcommand{\EtwoPriorAffineFourSmallPrior}{4.4}  % E2: 4 marks, 2 nm noise, 64-run prior + affine fit [nm]
\newcommand{\EtwoOursNineSmallPrior}{3.6}  % E2: 9 marks, 2 nm noise, 64-run prior, ours [nm]
\newcommand{\EtwoPriorAffineNineSmallPrior}{4.4}  % E2: 9 marks, 2 nm noise, 64-run prior + affine fit [nm]
\newcommand{\EtwoOursSixteenSmallPrior}{1.7}  % E2: 16 marks, 2 nm noise, 64-run prior, ours [nm]
\newcommand{\EtwoPriorAffineSixteenSmallPrior}{3.9}  % E2: 16 marks, 2 nm noise, 64-run prior + affine fit [nm]
\newcommand{\EtwoPolyTwoNine}{7.0}  % E2: 9 marks, 2 nm noise, 2nd-order correction [nm]
\newcommand{\EtwoPolyThreeSixteen}{2.2}  % E2: 16 marks, 2 nm noise, 3rd-order correction [nm]
\newcommand{\EtwoAffineSixteen}{5.8}  % E2: 16 marks, 2 nm noise, affine [nm]
\newcommand{\EthreeOodDTOurs}{0.96}  % E3 ood_dT_head_gt40: R2(d_max), ours
\newcommand{\EthreeOodDTBaseMin}{0.69}  % E3 ood_dT_head_gt40: R2(d_max), worst learned baseline
\newcommand{\EthreeOodDTBaseMax}{0.82}  % E3 ood_dT_head_gt40: R2(d_max), best learned baseline
\newcommand{\EthreeOodLOurs}{0.96}  % E3 ood_L_die_gt14: R2(d_max), ours
\newcommand{\EthreeOodLBaseMin}{0.66}  % E3 ood_L_die_gt14: R2(d_max), worst learned baseline
\newcommand{\EthreeOodLBaseMax}{0.92}  % E3 ood_L_die_gt14: R2(d_max), best learned baseline
\newcommand{\EthreeNoiseFiveOurs}{0.99}  % E3 train_noise_5nm: R2(d_max), ours
\newcommand{\EthreeNoiseFiveBaseMin}{0.83}  % E3 train_noise_5nm: R2(d_max), worst learned baseline
\newcommand{\EthreeNoiseFiveBaseMax}{0.90}  % E3 train_noise_5nm: R2(d_max), best learned baseline
\newcommand{\PinnFieldPhys}{92}  % E4: PINN field RMSE at N=64, with physics loss [nm]
\newcommand{\PinnFieldNoPhys}{88}  % E4: PINN field RMSE at N=64, physics loss removed [nm]
\newcommand{\PinnDmaxSixtyFour}{0.83}  % E4: PINN R2(d_max) at N=64
\newcommand{\NParams}{16}  % swept process/geometry/material parameters
\newcommand{\NTrainRuns}{1024}  % FEM training runs (Sobol)
\newcommand{\NHoldoutRuns}{128}  % FEM holdout runs (Latin hypercube)
\newcommand{\NPads}{21}  % pads per die edge (n x n grid)
\newcommand{\MeshCoarse}{12}  % elements along the die edge in the datasets
\newcommand{\RunTimeMedianS}{23}  % median wall time per FEM run [s]
\newcommand{\MeshFine}{28}  % finest mesh in the convergence study
\newcommand{\ConvCoarsePct}{0.2}  % worst m_x/d_max change, coarse vs finest mesh [%]
\newcommand{\DofsCoarse}{99,858}  % displacement DOFs, die + wafer, coarse mesh
\newcommand{\BondLayerEffectPct}{0.3}  % largest d_max change when the bond layers get Si properties [%]
\newcommand{\SiE}{130}  % Si Young's modulus [GPa]
\newcommand{\SiNu}{0.28}  % Si Poisson ratio
\newcommand{\SiAlpha}{2.6}  % Si thermal expansion [ppm/K]
\newcommand{\CuE}{120}  % Cu Young's modulus [GPa]
\newcommand{\CuAlpha}{16.7}  % Cu thermal expansion [ppm/K]
\newcommand{\Tref}{20}  % reference (pad patterning) temperature [C]
\newcommand{\WaferPatchFactor}{2}  % wafer patch edge / die edge
\newcommand{\NzDie}{4}  % elements through the die thickness
\newcommand{\NzWafer}{4}  % elements through the wafer thickness
\newcommand{\NRing}{6}  % elements across the wafer margin
\newcommand{\NTimeSteps}{20}  % thermal time steps per run
\newcommand{\TFirstUs}{1}  % first thermal time step [us]
\newcommand{\VerifTolPct}{1}  % tolerance of the analytic checks in tests/test_verification.py [%]
